\documentclass[main.tex]{subfiles}
\begin{document}

\begin{abstract}
  In recent years, significant advances have been made in numerical debugging
  and static analysis tools. Recent debugging tools significantly reduce
  overhead using double-double oracles, while recent static analysis tools
  accurately bound rounding errors using condition numbers. However, these new
  techniques also have downsides: double-double oracles can miss floating-point
  errors due to correlated errors, while condition number formalisms cannot
  cleanly handle over- and underflow. We show that combining both techniques
  avoids the downsides of each and implement this combination in \name. \name
  uses double-double computations but not as an oracle; instead, it detects
  erroneous operations using condition numbers, which minimizes overhead while
  achieving precision and recall similar to more expensive methods. \name also
  introduces a separate oracle to detect harmful over- and underflows more
  accurately; we implement this oracle in a new performant number
  representation. As a result, \name achieves a precision of \filled{80.0\%} and
  a recall of \filled{96.1\%} on a collection of 500 difficult numeric
  benchmarks. An ablation study shows that \name combines the strengths of both
  double-double oracles and condition numbers: it is more accurate than
  double-double oracles yet dramatically faster than methods based on
  arbitrary-precision condition number computations.
\end{abstract}

\end{document}
